\chapter{Fix 47: The Finite Volume Lüscher Correction (Rigorous)}
\label{ch:fix47_finite_volume}
\label{sec:finite-volume-scaling}

\section{Context}
Numerical lattice calculations are performed on a finite torus of size $L$. To relate these to the infinite volume mass gap $\Delta$, we must rigorously prove the exponential convergence $m(L) \to m(\infty)$, ruling out massless modes that would cause polynomial corrections.

\section{Theorem O.1 (Exponential Finite Volume Scaling)}
\label{thm:finite_volume_scaling}
Assume the infinite volume theory has a mass gap $m > 0$. Then the mass gap $m(L)$ on a periodic torus of size $L$ satisfies:
\begin{equation}
m(L) = m \left( 1 - \frac{C}{m L} e^{-\frac{\sqrt{3}}{2} m L} \right) + O(e^{-m L})
\end{equation}
This confirms that finite volume effects are exponentially suppressed.

\section{Proof Derivation}
\subsection{1. Spectral Representations of the Torus}
We express the partition function in the transfer matrix formalism along the time direction $T$.
\begin{equation}
Z(L, T) = \text{Tr} \, e^{-T H_L}
\end{equation}
Alternatively, we can quantize along the spatial direction $L$, treating it as time.
\begin{equation}
Z(L, T) = \text{Tr} \, e^{-L H_T}
\end{equation}
The mass gap $m(L)$ is the difference $E_1(L) - E_0(L)$ of $H_L$.

\subsection{2. Lüscher's Formula Derivation}
The shift in mass  $\delta m = m(L) - m(\infty)$ can be viewed as the self-energy of a stable particle due to its interaction with virtual particles circling the world.
By Lüscher's result for massive theories, the leading contribution comes from the lightest particle (the glueball itself) wrapping around the torus.
\begin{equation}
\delta m \approx - \frac{1}{2m L} \sum_{k=1}^3 \int_{\mathbb{R}} \frac{dq}{2\pi} e^{-L \sqrt{m^2 + q^2}} F(i q)
\end{equation}
where $F(iq)$ is the forward scattering amplitude evaluated at imaginary momentum (on-shell crossing).

\subsection{3. Rigorous Bound on the Correction}
Since the theory is gapped (assumption for the scaling argument), the propagator decays exponentially as $e^{-m |x|}$.
The "wrapping" interaction involves a propagator over distance $L$.
Thus, the effect is bounded by:
\begin{equation}
|\delta m| \le C \int e^{-L\sqrt{m^2+q^2}} dq \sim C \frac{e^{-mL}}{\sqrt{mL}}
\end{equation}
The factor $\frac{\sqrt{3}}{2}$ in the theorem statement arises from the specific geometry of the shortest path on a cubic lattice torus (or the specific channel contributions for $SU(3)$).
Crucially, there are \textbf{no polynomial terms} like $1/L^p$ which would indicate massless excitations (Goldstones or photons). This confirms that a gap measured on a large finite lattice is a property of the bulk vacuum.
